\section{Related Work}
\label{sec:relatedwork}
% EDIT (2026-10-29): moved from Section 2 to after the discussion and rewritten in plainer language,
% as asked in review. Same citations; the terms "gaming" and "reward hacking" are now defined before
% they are used (gaming in the introduction, reward hacking here).

\paragraph{Code models and their benchmarks.}
Code-generation models are usually scored by whether their output is functionally correct on small,
self-contained problems---HumanEval and its successors \cite{chen2021codex}---or, more recently, by
whether they can resolve an issue in a repository \cite{jimenez2024swebench}. Neither rewards
performance: a program that passes the tests scores the same whether it runs in one second or ten.
Our setting inverts this. Correctness is a gate that a candidate must pass, and the score is the
measured execution time of a real application.

\paragraph{LLMs for performance optimization.}
Work that uses LLMs to make code faster falls into three groups. The first learns from human
optimizations: PIE collects edits that made competitive-programming submissions faster and trains
models to reproduce them, reporting large speedups in simulated CPU time \cite{shypula2024pie}. The
second works on the compiler rather than the source: Cummins et al.\ train models to choose
optimization flags for LLVM code \cite{cummins2023llmopt} and release models trained on compiler
intermediate representations \cite{cummins2024llmcompiler}. The third writes GPU kernels directly.
KernelBench asks models to write CUDA kernels that match PyTorch reference modules and finds that
frontier models beat the PyTorch baseline in fewer than $20\%$ of tasks \cite{ouyang2025kernelbench};
ParEval tests parallel-code generation across several execution models and finds a large gap between
how well models write sequential and parallel code \cite{nichols2024parallel}. In all three groups the
program being optimized is a standalone kernel or a short program. In our study it is one file inside
a production application---AutoDock-GPU or Cholla---which keeps its build system, its callers and its
own scientific validation.

\paragraph{Profiling data as prompt context.}
The closest work to ours is KernelPro \cite{gai2026kernelpro}. It feeds Nsight Compute, SASS and
Nsight Systems output through a set of small tools that turn raw metrics into guidance in natural
language, chooses which tools to run with a roofline-based classifier, and reports a $125\%$ higher
speedup for the guidance than for the raw metrics. Both of its arms carry profiling data; the
comparison is interpreted against raw. We ask a narrower and, we argue, prior question: does profiling
data help \emph{at all}, compared with the same prompt without it? E4 answers it by removing the data
while holding every other part of the prompt fixed, and finds that the answer depends on the workload:
no net effect on single-kernel benchmarks, where the data moves individual results in both directions,
but a consistent gain on a large application where the bottleneck is not apparent from the source. E5
then evaluates KernelPro's question on our benchmarks and finds that interpretation widens the spread
of outcomes rather than shifting its centre.

\paragraph{Autotuning and search.}
Classical autotuners make code faster by trying many versions and measuring each: OpenTuner for
general programs \cite{ansel2014opentuner}, and AutoTVM \cite{chen2018autotvm} and Ansor
\cite{zheng2020ansor} for tensor programs, where a learned cost model predicts which versions are
worth measuring. These systems search a space that a human designed in advance and need many
measurements; an LLM instead proposes a small number of edits that can change the program in any way.
The two are complementary. Our build--verify--measure loop is the autotuner's evaluation harness
wrapped around a model that proposes the versions, and the profile summaries in the prompt play the
role the cost model plays in search: they tell the proposer where the time goes. Roofline analysis
\cite{williams2009roofline} underlies both the Nsight Compute reports we summarize and the bottleneck
classification used by prior work.

\paragraph{Repair loops.}
Letting a model fix its own output from compiler and test feedback is now standard practice
\cite{chen2024selfdebug}. Our loop is a narrow instance: the feedback is a build log or an
application-level correctness gate, and the number of repair iterations is bounded. What we add is
practical rather than a new technique: a list of cases in which the loop cannot work because the
information the model would need is missing from the feedback, for example a linker error whose
explanatory line is filtered out of the summary, or an application that aborts with a bare assertion
and discards the CUDA error (Section~\ref{sec:pitfalls}).

\paragraph{Benchmark gaming and reward hacking.}
\emph{Reward hacking} is the general name for what happens when a system that is optimized against a
measurable objective finds a way to score well on the measure without doing the task the measure was
meant to capture. Benchmark gaming (Section~\ref{sec:intro}) is its code-generation instance. Recent
benchmarks study it by exposing evaluation harnesses that can be exploited: EvilGenie places coding
agents in an environment where test files can be edited or answers hardcoded, and finds that current
proprietary agents do so, in ways an LLM judge detects but held-out tests do not reliably
\cite{gabor2025evilgenie}. In the performance setting specifically, KernelBench-Verified audits
LLM-generated GPU kernels and finds identity functions reported as $374\times$ speedups, fast paths
that fire only for the test tensor's exact dimensions, and work skipped depending on input values;
scored against hidden input distributions and an LLM-judge audit, the best model drops from
$1.43\times$ to $0.88\times$, with $11.6\%$ of problems relying on such exploits
\cite{zhang2026kernelbenchverified}. Our contribution is therefore not the observation but its setting
and its remedy: in an HPC benchmark, deleting the measured computation preserves the printed output and
the benchmark's own PASS check while producing three-order-of-magnitude ``speedups'' (E2); a
prompt-level semantic-preservation constraint, which we measure, stops it for both models we tested;
and the detection rule we use in its place---an implausibility threshold followed by reading the
diff---is cheap enough to run on every candidate.

\paragraph{Repository-level optimization.}
A parallel line of work optimizes whole repositories rather than kernels. GSO scores agents against the
speedup an expert commit achieved across ten codebases \cite{shetty2025gso}, and PerfAgent adds a
profiler and a verifier to the agent's loop and more than doubles the rate of patches that match the
expert's \cite{deng2026perfagent}. Its motivation matches ours almost exactly---agents miss bottlenecks
hidden behind abstraction layers and native extensions---but its targets are CPU-bound Python
repositories scored against an expert's commit. Ours are GPU kernels inside production scientific
applications, scored against the application's own scientific output, which is a correctness standard
an expert diff cannot provide.

\paragraph{HPC benchmark suites.}
We use HeCBench \cite{jin2023hecbench}, which provides many small CUDA kernels with built-in
verification, for the benchmark experiments, and two production applications for the end-to-end study.
The contrast between the two is one of our main findings: conclusions drawn on single-kernel
benchmarks did not transfer to applications, either for the value of profiling data or for the prompt
variant that works best.
